\chapter{INTRODUCTION}
\pagebreak

\begin{center}
{\LARGE\textbf{INTRODUCTION}}
\end{center}

\textcolor{red}{DISCLAIMER: The outline given here is nothing but a suggestion. Please consult your advisor before using them.}

Saudi Arabia is using more Internet of Things (IoT) devices for its smart cities and projects under Vision 2030. However, these devices are often not secure. A major example is the 2016 Mirai botnet attack, which took control of over 600,000 IoT devices worldwide to cause major internet outages [2]. This shows that insecure IoT devices are a real and serious problem that can also affect companies and infrastructure in Saudi Arabia.

Security tools called honeypots can be used to attract and study attackers. Other researchers have also worked on using artificial intelligence (AI) to understand attacks. But these are usually separate solutions. Our project, TrapPot, combines these ideas. It is a smart honeypot system that uses AI to not only detect attacks on IoT devices but also to automatically identify the type of attack and create useful reports. This will help make Saudi Arabia's growing digital environment more secure. \cite{203628}.

% ============================================================
\section{Challenges}
\label{sec:challenges}
% ============================================================


Several significant challenges are inherent to developing an IoT honeypot system. First, the technical complexity of the subject is high, requiring expertise in cybersecurity, networking, and artificial intelligence, which can be demanding for undergraduate IT students. Second, the scarcity of specialized resources presents a hurdle, as this is a niche field with limited tutorials and case studies compared to more common IT topics. Third, there is a critical security risk; if the honeypot is not designed and isolated correctly, it could itself become a vulnerability and be used as an entry point into our own systems. Finally, the data management aspect is a major challenge, as the system needs to efficiently collect, store, and process potentially large volumes of attack logs to feed the AI module.

Our project directly addresses the technical and data management challenges through careful planning and the use of established tools. We are mitigating the security risk by designing the system with strict network isolation in mind. The scarcity of resources remains an external challenge we must navigate through focused research.

% ============================================================
\section{Study Scope}
\label{sec:study_scope}
% ============================================================

This project focuses on developing an AI-powered honeypot system that emulates Telnet and SSH services to capture IoT-targeted attacks in an isolated network environment. The scope is limited to implementing Random Forest classification for four attack types (brute-force, scanning, malware download, and command execution), collecting data over a 30-day period, and achieving minimum 85% classification accuracy. The system operates as a single-node proof-of-concept with complete network isolation to prevent production system compromise. Advanced protocols (HTTP, MQTT, CoAP), deep learning methods, distributed deployments, and production-ready features are excluded from this phase and reserved for future work.
Random Forest has already shown high accuracy and low latency for classifying honeypot‑captured IoT botnet attacks, making it a suitable choice for TrapPot’s classifier \cite{lee2021classification}. Telnet and SSH services are also the primary targets in Mirai‑like IoT botnet campaigns and SSH honeypot studies, which motivates our focus on these protocols and the selected attack types \cite{herwig2021analyzing,zhang2023identification}.

% ============================================================
\section{Goals and Objectives}
\label{sec:goals_and_objectives}
% ============================================================

The project's main goals are to deploy an isolated IoT honeypot emulating Telnet and SSH services, automatically classify captured attacks using machine learning, achieve classification accuracy exceeding 85% , generate real-time threat intelligence reports, and ensure complete network isolation to prevent compromise of production systems.

% ============================================================
\section{Report Arrangement}
\label{sec:report_arrangement}
% ============================================================

This report is structured to guide the reader through the conception, design, and proposed implementation of the TrapPot system, as outlined below:

Chapter 1: Introduction - Establishes the core problem of IoT security, the project's motivation, and its primary objectives. It defines the project's boundaries and presents this overall report structure.

Chapter 2: Literature Review - Surveys the existing academic and technical landscape, covering foundational work on IoT honeypots, machine learning for attack classification, and related systems to position TrapPot within the field.

Chapter 3: Methodology - Details the architectural design and operational plan for TrapPot. It explains the system components, the data processing pipeline from detection to classification, and the proposed evaluation metrics.

Chapter 4: Implementation and Testing - (For Future Work) This chapter will document the practical development, specific tools, coding frameworks, and the comprehensive testing strategy used to validate the system.

Chapter 5: Results and Discussion - (For Future Work) This chapter will present and analyze the outcomes of the testing phase, interpreting the system's performance and effectiveness.

Chapter 6: Conclusion and Future Work - (For Future Work) This chapter will summarize the project's contributions, draw final conclusions, and propose directions for further development and research.
